\section{Introduction to Linear Higher Order Equations}

\noindent
An $n$th order differential equation is said to be {\color{blue}\it linear} if it
can be written in the form
\begin{equation}\label{eq:9.1.1}
y^{(n)}+p_1(x)y^{(n-1)}+\cdots+p_n(x)y=f(x).
\end{equation}
We considered equations of this form with $n=1$ in
Section~2.1 and with $n=2$ in
Chapter~5. In this chapter $n$ is an arbitrary
positive integer.

In this section we sketch the general theory of linear $n$th order
equations. Since this theory has already been discussed for $n=2$ in
Sections~5.1\; and 5.3, we'll omit
proofs.

For convenience, we consider linear differential equations written
as
\begin{equation}\label{eq:9.1.2}
P_0(x)y^{(n)}+P_1(x)y^{(n-1)}+\cdots+P_n(x)y=F(x),
\end{equation}
which can be rewritten as \eqref{eq:9.1.1} on any interval on
which $P_0$ has no zeros, with $p_1=P_1/P_0$, \dots, $p_n=P_n/P_0$ and
$f=F/P_0$. For simplicity, throughout this chapter we'll
abbreviate the left side of \eqref{eq:9.1.2} by $Ly$; that is,
$$
Ly=P_0y^{(n)}+P_1y^{(n-1)}+\cdots+P_ny.
$$
We say that the equation $Ly=F$ is {\color{blue}\it normal\/} on $(a,b)$ if
$P_0$, $P_1$, \dots, $P_n$ and $F$ are continuous on $(a,b)$ and $P_0$ has
no zeros on $(a,b)$. If this is so then $Ly=F$ can be
written as \eqref{eq:9.1.1} with $p_1$, \dots, $p_n$ and $f$ continuous on
$(a,b)$.

The next theorem is analogous to
Theorem~\ref{thmtype:5.3.1}.

\begin{theorem}\color{blue}\label{thmtype:9.1.1}
Suppose   $Ly=F$ is normal on $(a,b)$, let $x_0$ be a point in
$(a,b),$
and let $k_0$, $k_1$, \dots, $k_{n-1}$ be arbitrary real numbers$.$ Then
the initial value problem
$$
Ly=F, \quad  y(x_0)=k_0,\quad y'(x_0)=k_1,\dots,\quad y^{(n-1)}(x_0)=k_{n-1}
$$
has a unique solution on $(a,b).$
\end{theorem}

\boxit{Homogeneous Equations}

\noindent
Eqn.~\eqref{eq:9.1.2} is said to be {\color{blue}\it homogeneous} if $F\equiv0$ and
{\color{blue}\it nonhomogeneous} otherwise. Since $y\equiv0$ is obviously a
solution of $Ly=0$, we call it the {\color{blue}\it trivial} solution. Any other
solution is {\color{blue}\it nontrivial}.

If $y_1$, $y_2$, \dots, $y_n$ are defined on $(a,b)$ and
$c_1$, $c_2$, \dots, $c_n$ are constants, then
\begin{equation}\label{eq:9.1.3}
y=c_1y_1+c_2y_2+\cdots+c_ny_n
\end{equation}
is a {\color{blue}\it linear combination\/} of $\{y_1,y_2\dots,y_n\}$. It's
easy
 to show that if $y_1$, $y_2$, \dots, $y_n$ are solutions of
$Ly=0$ on $(a,b)$, then so is any linear combination of
$\{y_1,y_2,\dots,y_n\}$. (See the proof of Theorem~\ref{thmtype:5.1.2}.) We
say that $\{y_1,y_2,\dots,y_n\}$ is a {\color{blue}\it fundamental set of
solutions of $Ly=0$ on\/} $(a,b)$ if every solution of $Ly=0$ on
$(a,b)$ can be written as a linear combination of
$\{y_1,y_2,\dots,y_n\}$, as in \eqref{eq:9.1.3}. In this case we say that
\eqref{eq:9.1.3} is the {\color{blue}\it general solution of $Ly=0$ on
$(a,b)$\/}.

It can be shown (Exercises~\ref{exer:9.1.14} and \ref{exer:9.1.15}) that if
the equation $Ly=0$ is normal on $(a,b)$ then it has
infinitely many fundamental  sets of solutions  on $(a,b)$.
The next definition will help to identify fundamental sets of
solutions of $Ly=0$.

We say that $\{y_1,y_2,\dots,y_n\}$ is {\color{blue}\it linearly independent\/} on
$(a,b)$ if the only constants $c_1$, $c_2$, \dots, $c_n$ such that
\begin{equation}\label{eq:9.1.4}
c_1y_1(x)+c_2y_2(x)+\cdots+c_ny_n(x)=0,\quad a<x<b,
\end{equation}
are $c_1=c_2=\cdots=c_n=0$. If \eqref{eq:9.1.4} holds for some set of
constants $c_1$, $c_2$, \dots, $c_n$ that are not all zero, then
$\{y_1,y_2,\dots,y_n\}$ is {\color{blue}\it linearly dependent\/} on $(a,b)$

The next theorem is analogous to
Theorem~\ref{thmtype:5.1.3}.

\begin{theorem}\color{blue}\label{thmtype:9.1.2}
If  $Ly=0$ is normal on $(a,b)$, then a set
$\{y_1,y_2,\dots,y_n\}$ of $n$ solutions of $Ly=0$ on $(a,b)$ is a
fundamental set if and only if it's linearly independent on $(a,b)$.
\end{theorem}

\begin{example}\label{example:9.1.1}\rm
The equation
\begin{equation}\label{eq:9.1.5}
x^3y'''-x^2y''-2xy'+6y=0
\end{equation}
is normal and has the solutions $y_1=x^2$, $y_2=x^3$, and $y_3=1/x$ on
$(-\infty,0)$ and $(0,\infty)$.  Show that $\{y_1,y_2,y_3\}$
is linearly independent on $(-\infty, 0)$ and $(0,\infty)$. Then find
the general solution of \eqref{eq:9.1.5} on $(-\infty, 0)$ and
$(0,\infty)$.
\end{example}

\solution
Suppose
\begin{equation} \label{eq:9.1.6}
c_1x^2+c_2x^3+{c_3\over x}=0
\end{equation}
on $(0,\infty)$. We must show that $c_1=c_2=c_3=0$. Differentiating
\eqref{eq:9.1.6} twice yields  the system
\begin{equation} \label{eq:9.1.7}
\begin{array}{rcl}
c_1x^2+c_2x^3+\dst{c_3\over x}&=&0\\[2\jot]
2c_1x+3c_2x^2-\dst{c_3\over x^2}&=&0\\[2\jot]
2c_1+6c_2x+\dst{2c_3\over x^3}&=&0.
\end{array}
\end{equation}
If \eqref{eq:9.1.7} holds for all $x$ in $(0,\infty)$, then it
certainly  holds at $x=1$;  therefore,
\begin{equation} \label{eq:9.1.8}
\begin{array}{rcl}
\phantom{2}c_1+\phantom{3}c_2+\phantom{2}c_3&=&0\\
2c_1+3c_2-\phantom{2}c_3&=&0\\
2c_1+6c_2+2c_3&=&0.
\end{array}
\end{equation}
By solving this system directly, you can verify that it has only the
trivial solution $c_1=c_2=c_3=0$;     however, for our purposes it's more
useful to recall from linear algebra that a homogeneous linear system
of $n$ equations in $n$ unknowns has only the trivial solution if its
determinant is nonzero. Since the determinant of \eqref{eq:9.1.8} is
$$
\left|\begin{array}{rrr}1&1&1\\2&3&-1\\2&6&2\end{array}\right|=
\left|\begin{array}{rrr}1&0&0\\2&1&-3\\2&4&0\end{array}\right|=12,
$$
it follows that \eqref{eq:9.1.8} has only the trivial solution, so
$\{y_1,y_2,y_3\}$ is linearly independent on $(0,\infty)$. Now
Theorem~\ref{thmtype:9.1.2} implies that
$$
y=c_1x^2+c_2x^3+{c_3\over x}
$$
is the general solution of \eqref{eq:9.1.5} on $(0,\infty)$.
To see that this is also true
 on $(-\infty,0)$,  assume that \eqref{eq:9.1.6} holds on
$(-\infty,0)$. Setting $x=-1$ in \eqref{eq:9.1.7} yields
\begin{eqnarray*}
\phantom{-2}c_1-\phantom{3}c_2-\phantom{2}c_3&=&0\\
-2c_1+3c_2-\phantom{2}c_3&=&0\\
\phantom{-}2c_1-6c_2-2c_3&=&0.
\end{eqnarray*}
Since the determinant of this system is
$$
\left|\begin{array}{rrr}1&-1&-1\\-2&3&-1\\2&-6&-2\end{array}\right|=
\left|\begin{array}{rrr}1&0&0\\-2&1&-3\\2&-4&0\end{array}\right|=-12,
$$
it follows that $c_1=c_2=c_3=0$; that is, $\{y_1,y_2,y_3\}$
is linearly independent on $(-\infty,0)$.

\begin{example}\label{example:9.1.2}\rm
The equation
\begin{equation}\label{eq:9.1.9}
y^{(4)}+y'''-7y''-y'+6y=0
\end{equation}
is normal and has the solutions $y_1=e^x$, $y_2=e^{-x}$, $y_3=e^{2x}$,
and $y_4=e^{-3x}$ on $(-\infty,\infty)$. (Verify.) Show that
$\{y_1,y_2,y_3,y_4\}$ is linearly independent on $(-\infty,\infty)$.
Then find the general solution of \eqref{eq:9.1.9}.
\end{example}

\solution
Suppose   $c_1$, $c_2$, $c_3$, and $c_4$
are constants such that
\begin{equation} \label{eq:9.1.10}
c_1e^x+c_2e^{-x}+c_3e^{2x}+c_4e^{-3x}=0
\end{equation}
for all $x$. We must show that $c_1=c_2=c_3=c_4=0$. Differentiating
\eqref{eq:9.1.10} three times yields the system
\begin{equation} \label{eq:9.1.11}
\begin{array}{rcl}
c_1e^x+c_2e^{-x}+\phantom{2}c_3e^{2x}+\phantom{27}c_4e^{-3x}&=&0\\
c_1e^x-c_2e^{-x}+2c_3e^{2x}-\phantom{2}3c_4e^{-3x}&=&0\\
c_1e^x+c_2e^{-x}+4c_3e^{2x}+\phantom{2}9c_4e^{-3x}&=&0\\
c_1e^x-c_2e^{-x}+8c_3e^{2x}-27c_4e^{-3x}&=&0.
\end{array}
\end{equation}
If \eqref{eq:9.1.11} holds for all $x$, then it certainly holds for $x=0$.
Therefore
$$
\begin{array}{rcl}
c_1+c_2+\phantom{2}c_3+\phantom{27}c_4&=&0\\
c_1-c_2+2c_3-\phantom{2}3c_4&=&0\\
c_1+c_2+4c_3+\phantom{2}9c_4&=&0\\
c_1-c_2+8c_3-27c_4&=&0.
\end{array}
$$
The determinant of this system is
\begin{equation} \label{eq:9.1.12}
\begin{array}{rcl}
\left|\begin{array}{rrrr}1&1&1&1\\1&-1&2&-3\\1&1&4&9\\
1&-1&8&-27\end{array}\right|&=&
\left|\begin{array}{rrrr}1&1&1&1\\0&-2&1&-4\\0&0&3&8\\
0&-2&7&-28\end{array}\right|
=\left|\begin{array}{rrr}-2&1&-4\\0&3&8\\
-2&7&-28\end{array}\right|\\[8\jot]&=&
\left|\begin{array}{rrr}-2&1&-4\\0&3&8
\\0&6&-24\end{array}\right|=
-2\left|\begin{array}{rr}3&8\\6&-24\end{array}\right|=240,
\end{array}
\end{equation}
so the system has only the trivial solution $c_1=c_2=c_3=c_4=0$.
Now Theorem~\ref{thmtype:9.1.2} implies that
$$
y=c_1e^x+c_2e^{-x}+c_3e^{2x}+c_4e^{-3x}
$$
is the general solution of \eqref{eq:9.1.9}.

\boxit{The Wronskian}

\noindent
We can use the method used in Examples~\ref{example:9.1.1} and
\ref{example:9.1.2} to test $n$ solutions $\{y_1,y_2,\dots,y_n\}$ of any
$n$th order equation $Ly=0$ for linear independence on an interval
$(a,b)$ on which the equation is normal. Thus, if $c_1$, $c_2$ ,\dots,
$c_n$ are constants such that
$$
c_1y_1+c_2y_2+\cdots+c_ny_n=0,\quad a<x<b,
$$
then differentiating $n-1$ times leads to the $n\times n$
system of equations
\begin{equation} \label{eq:9.1.13}
\begin{array}{rcl}
c_1y_1(x)+c_2y_2(x)+&\cdots&+c_ny_n(x)=0\\
c_1y_1'(x)+c_2y_2'(x)+&\cdots&+c_ny_n'(x)=0\\
\phantom{c_1y_1^{(n)}(x)+c_2y_2^{(n-1)}(x)}&\vdots&
\phantom{\cdots+c_ny_n^{(n-1)(x)}=q}\\
c_1y_1^{(n-1)}(x)+c_2y_2^{(n-1)}(x)+&\cdots&+c_ny_n^{(n-1)}(x)
=0
\end{array}
\end{equation}
for $c_1$, $c_2$, \dots, $c_n$. For a fixed $x$,  the determinant of this
system is
$$
W(x)=\left|\begin{array}{cccc}
y_1(x)&y_2(x)&\cdots&y_n(x)\\[2\jot]
y'_1(x)&y'_2(x)&\cdots&y_n'(x)\\[2\jot]
\vdots&\vdots&\ddots&\vdots\\[2\jot]
y_1^{(n-1)}(x)&y_2^{(n-1)}(x)&\cdots&y_n^{(n-1)}(x)
\end{array}\right|.
$$
We call this determinant the
\href{https://mathshistory.st-andrews.ac.uk/Biographies/Wronski/}
{\color{blue}\it Wronskian\/} of
$\{y_1,y_2,\dots,y_n\}$. If $W(x)\ne0$ for some $x$ in $(a,b)$ then
 the system \eqref{eq:9.1.13} has only the trivial solution
$c_1=c_2=\cdots=c_n=0$, and Theorem~\ref{thmtype:9.1.2} implies that
$$
y=c_1y_1+c_2y_2+\cdots+c_ny_n
$$
is the general solution of $Ly=0$ on $(a,b)$.

The next theorem generalizes Theorem~\ref{thmtype:5.1.4}.
The proof is sketched in (Exercises~\ref{exer:9.1.17}--\ref{exer:9.1.20}).

\begin{theorem}\color{blue}\label{thmtype:9.1.3}
 Suppose the   homogeneous linear $n$th order equation
\begin{equation}\label{eq:9.1.14}
P_0(x)y^{(n)}+P_1(x)y^{n-1}+\cdots+P_n(x)y=0
\end{equation}
is normal on $(a,b),$ let $y_1,$ $y_2,$ \dots, $y_n$ be solutions of
\eqref{eq:9.1.14} on $(a,b),$ and let $x_0$ be in $(a,b)$. Then the
Wronskian of $\{y_1,y_2,\dots,y_n\}$ is given by
\begin{equation}  \label{eq:9.1.15}
W(x)=W(x_0)\exp\left\{-\int^x_{x_0}{P_1(t)\over P_0(t)}\,
dt\right\},\quad a<x<b.
\end{equation}
Therefore$,$  either $W$ has no zeros in  $(a,b)$ or $W\equiv0$
on  $(a,b).$
\end{theorem}

Formula \eqref{eq:9.1.15} is
 \href{https://mathshistory.st-andrews.ac.uk/Biographies/Abel/}
{\color{blue}\it Abel's formula\/}.

The next theorem is analogous to
Theorem~\ref{thmtype:5.1.6}..

\begin{theorem}\color{blue}\label{thmtype:9.1.4}
Suppose $Ly=0$ is normal on $(a,b)$ and let $y_1$, $y_2$, \dots, $y_n$
be $n$ solutions of $Ly=0$ on $(a,b)$. Then the following statements
are equivalent$;$  that is$,$ they are either all true or all false$:$
\begin{alist}
\item % (a)
The general solution of $Ly=0$ on $(a,b)$ is
$y=c_1y_1+c_2y_2+\cdots+c_ny_n.$
\item % (b)
 $\{y_1,y_2,\dots,y_n\}$ is a fundamental set of solutions of
$Ly=0$ on $(a,b).$
\item % (c)
$\{y_1,y_2,\dots,y_n\}$ is linearly independent on $(a,b).$
\item % (d)
The Wronskian of $\{y_1,y_2,\dots,y_n\}$ is nonzero at some point in
$(a,b).$
\item % (e)
The Wronskian of $\{y_1,y_2,\dots,y_n\}$ is nonzero at all points in
$(a,b).$
\end{alist}
\end{theorem}

\begin{example}\label{example:9.1.3} \rm
In Example~\ref{example:9.1.1} we saw that the solutions $y_1=x^2$,
$y_2=x^3$, and $y_3=1/x$ of
$$
x^3y'''-x^2y''-2xy'+6y=0
$$
are linearly independent on $(-\infty,0)$ and $(0,\infty)$. Calculate
the Wronskian  of $\{y_1,y_2,y_3\}$.
\end{example}

\solution
If $x\ne0$,  then
$$
W(x)=\left|\begin{array}{ccc}
x^2&\phantom{2}x^3&\phantom{-}\dst{1\over x}\\[2\jot]
2x&3x^2\phantom{^2}&\dst-\dst{1\over x^2}\\[2\jot]
2&6x\phantom{x^2}&\phantom{-}\dst{2\over x^3}\end{array}
\right|
=2x^3\left|\begin{array}{ccc}
1&\phantom{3}x&\phantom{-}\dst{1\over x^3}\\[2\jot]2&3x&-\dst{1\over
x^3}
\\[2\jot]1&3x&\phantom{-}\dst{1\over x^3}\end{array}\right|,
$$
where we factored $x^2$, $x$, and $2$ out of the first, second, and
third  rows of
$W(x)$, respectively. Adding the second row of the last determinant
to the first and third rows yields
$$
W(x)=2x^3\left|\begin{array}{ccc}
3&4x&\phantom{-}0\\[2\jot]2&3x&-\dst{1\over
x^3}
\\[2\jot]3&6x&\phantom{-}0\end{array}\right|
=2x^3\left(1\over x^3\right)\left|\begin{array}{cc}3&4x\\3&6x
\end{array}\right|=12x.
$$
Therefore $W(x)\ne0$ on $(-\infty,0)$ and $(0,\infty)$.

\begin{example}\label{example:9.1.4} \rm In
Example~\ref{example:9.1.2} we saw that the solutions $y_1=e^x$,
$y_2=e^{-x}$,
$y_3=e^{2x}$, and $y_4=e^{-3x}$ of
$$
y^{(4)}+y'''-7y''-y'+6y=0
$$
are linearly independent on every open interval.
Calculate the Wronskian  of $\{y_1,y_2,y_3,y_4\}$.
\end{example}

\solution
For all $x$,
$$
W(x)=\left|\begin{array}{rrrr}
e^x&e^{-x}&e^{2x}&e^{-3x}\\
e^x&-e^{-x}&2e^{2x}&-3e^{-3x}\\
e^x&e^{-x}&4e^{2x}&9e^{-3x}\\
e^x&-e^{-x}&8e^{2x}&-27e^{-3x}
\end{array}\right|.
$$
Factoring the exponential common factor from each row yields
$$
W(x)=e^{-x}\left|\begin{array}{rrrr}1&1&1&1\\1&-1&2&-3\\1&1&4&9\\
1&-1&8&-27\end{array}\right|=240e^{-x},
$$
from \eqref{eq:9.1.12}.

\color{blue}
\remark{Under the assumptions of Theorem~\ref{thmtype:9.1.4}, it isn't
necessary to obtain a formula for $W(x)$. Just evaluate $W(x)$
at a convenient point in $(a,b)$, as we did in
Examples~\ref{example:9.1.1} and \ref{example:9.1.2}.}
\color{black}

\begin{theorem}\color{blue} \label{thmtype:9.1.5}
Suppose $c$ is in $(a,b)$ and $\alpha_{1},$ $\alpha_{2},$ \dots$,$
are real numbers, not all zero.
Under the assumptions of Theorem~\ref{thmtype:10.3.3}, suppose
$y_{1}$ and $y_{2}$ are  solutions  of \eqref{eq:5.1.35}   such that
\begin{equation} \label{eq:9.1.16}
\alpha y_{i}(c)+ y_{i}'(c)+\cdots +y_{i}^{(n-1)}(c)=0,\quad 1\le i\le n.
\end{equation}
Then $\{y_{1},y_{2},\dots y_{n}\}$ isn't  linearly independent on $(a,b).$
\end{theorem}

\proof Since $\alpha_{1}$, $\alpha_{2}$, \dots, $\alpha_{n}$
are not all zero,
\eqref{eq:9.1.14} implies that
$$
\left|\begin{array}{ccccccc}
y_{1}(c)&y_{1}'(c)&\cdots&y_{1}^{(n-1)}(c)\\
y_{2}(c)&y_{2}'(c)&\cdots&y_{2}^{(n-1)}(c)\\
\vdots&\vdots&\ddots&\vdots\\
y_{n}(c)&y_{n}'(c)&\cdots&y_{n}^{(n-1)}(c)\\
\end{array}\right|=0,
$$
so
$$
\left|\begin{array}{cccccc}
y_{1}(c)&y_{2}(c)&\cdots& y_{n}(c)\\
y_{1}'(c)&y_{2}'(c)&\cdots& y_{n}'(c)\\
\vdots&\vdots&\ddots&\vdots\\
y_{1}^{(n-1)}(c)&y_{2}^{(n-1)}(c)(c)&\cdots& y_{n}^{(n-1)}(c)(c)\\
\end{array}\right|=0
$$
\nopagebreak
and Theorem~\ref{thmtype:9.1.4} implies the stated conclusion.
\newpage

\boxit{General Solution of a Nonhomogeneous Equation}

\noindent
The next theorem is analogous to
Theorem~\ref{thmtype:5.3.2}.
It shows how to find the general solution of $Ly=F$ if we know a particular
solution of $Ly=F$ and a fundamental set of solutions of the {\color{blue}\it
complementary equation\/} $Ly=0$.

\begin{theorem}\color{blue}\label{thmtype:9.1.6}
Suppose $Ly=F$ is normal on $(a,b).$ Let $y_p$ be a particular
solution of $Ly=F$ on $(a,b),$ and let $\{y_1,y_2,\dots,y_n\}$ be a
fundamental set of solutions of the complementary equation $Ly=0$ on
$(a,b)$. Then $y$ is a solution of $Ly=F$ on $(a,b)$ if and only if
$$
y=y_p+c_1y_1+c_2y_2+\cdots+c_ny_n,
$$
where $c_1,c_2,\dots,c_n$  are constants.
\end{theorem}

The next theorem is analogous to Theorem~\ref{thmtype:5.3.2}.

\begin{theorem}\color{blue}$[$The Principle of
Superposition$]$\label{thmtype:9.1.7}
Suppose for each $i=1,$ $2,$ \dots, $r$, the function $y_{p_i}$ is a
particular solution of $Ly=F_i$ on  $(a,b).$ Then
$$
y_p=y_{p_1}+y_{p_2}+\cdots+y_{p_r}
$$
is a particular  solution of
$$
Ly=F_1(x)+F_2(x)+\cdots+F_r(x)
$$
on $(a,b).$
\end{theorem}

We'll apply  Theorems~\ref{thmtype:9.1.6} and \ref{thmtype:9.1.7}
throughout the rest of this chapter.

\exercises

\begin{exerciselist}

\item\label{exer:9.1.1}
Verify that the given function is the solution of the initial
value problem.

\part{a} $ x^3y'''-3x^2y''+6xy'-6y=-\dst{24\over x}, \quad  y(-1)=0$, $
y'(-1)=0, \quad  y''(-1)=0$;

\hspace*{15pt}$\dst{y=-6x-8x^2-3x^3+{1\over x}}$

\part{b} $\dst{ y'''-{1\over x}y''-y'+{1\over x}y={x^2-4\over x^4}, \quad
y(1)={3\over2},\quad y'(1)={1\over2}}$, $ y''(1)=1$;

\hspace*{15pt} $y=x+\dst{1\over2x}$

\part{c} $xy'''-y''-xy'+y=x^2, \quad  y(1)=2,\quad y'(1)=5,\quad
y''(1)=-1$;

\hspace*{15pt} $y=-x^2-2+2e^{(x-1)}-e^{-(x-1)}+4x$

\part{d} $4x^3y'''+4x^2y''-5xy'+2y=30x^2, \quad  y(1)=5,\quad
y'(1)=\dst{17\over2}$;

\hspace*{15pt} $y''(1)=\dst{63\over4};\quad y=2x^2\ln
x-x^{1/2}+2x^{-1/2}+4x^2$

\part{e} $x^4y^{(4)}-4x^3y'''+12x^2y''-24xy'+24y=6x^4, \quad  y(1)=-2$,

\hspace*{15pt} $y'(1)=-9, \quad  y''(1)=-27,\quad y'''(1)=-52$;

\hspace*{15pt} $y=x^4\ln x+x-2x^2+3x^3-4x^4$

\part{f}
$xy^{(4)}-y'''-4xy''+4y'=96x^2, \quad  y(1)=-5,\quad y'(1)=-24$

\hspace*{15pt} $y''(1)=-36; \quad   y'''(1)=-48;\quad y=9-12x+6x^2-8x^3$

\item\label{exer:9.1.2}
Solve the initial value problem
$$
x^3y'''-x^2 y''-2xy'+6y=0,
\quad   y(-1)=-4, \quad   y'(-1)=-14,\quad   y''(-1)=-20.
$$
\hint{See Example~$\ref{example:9.1.1}$.}

\item\label{exer:9.1.3}
Solve the initial value problem
$$
y^{(4)}+y'''-7y''-y'+6y=0,
\quad   y(0)=5,\quad   y'(0)=-6,\quad   y''(0)=10,\quad   y'''(0)-36.
$$
\hint{See Example~$\ref{example:9.1.2}$.}

\item\label{exer:9.1.4}
Find solutions $y_1$, $y_2$, \dots, $y_n$ of the equation $y^{(n)}=0$
that satisfy the initial conditions
$$
y_i^{(j)}(x_0)=\left\{\begin{array}{cl}
0,&j\ne i-1,\\[1\jot]
1,&j=i-1,\end{array}\right.\; 1\le i\le n.
$$

\item\label{exer:9.1.5}
\begin{alist}
\item %(a)
Verify that the function
$$
y=c_1x^3+c_2x^2+{c_3\over x}
$$
satisfies
$$
x^3 y'''-x^2y''-2xy'+6y=0
\eqno{\rm (A)}
$$
if $c_1$, $c_2$, and $c_3$ are constants.

\item %(b)
Use \part{a} to find solutions $y_1$, $y_2$, and $y_3$ of
(A) such that
$$
\begin{array}{rl}
y_1(1)&=1,\quad y_1'(1)=0,\quad  y_1''(1)=0  \\[1\jot]
y_2(1)&=0,\quad  y_2'(1)=1,\quad  y_2''(1)=0  \\[1\jot]
y_3(1)&=0,\quad  y_3'(1)=0,\quad  y_3''(1)=1.
\end{array}
$$

\item %(c)
Use \part{b} to find the solution of  (A) such that
$$
y(1)=k_0,\quad y'(1)=k_1,\quad   y''(1)=k_2.
$$
\end{alist}

\item\label{exer:9.1.6}
Verify that the given functions are solutions of the given equation, and
 show that they  form a fundamental set of solutions of the
equation on any interval on which the equation is normal.

\part{a} $y'''+y''-y'-y=0;  \quad\{e^x,\,e^{-x},\,xe^{-x}\}$

\part{b} $y'''-3y''+7y'-5y=0;  \quad\{e^x,\,e^x\cos2x,\,e^x\sin2x\}$.

\part{c} $xy'''-y''-xy'+y=0;  \quad \{e^x,\,e^{-x},\,x\}$

\part{d} $x^2y'''+2xy''-(x^2+2)y=0;  \quad\dst\{e^x/ x,\,e^{-x}/
x,\,1\}$

\part{e} $(x^2-2x+2)y'''-x^2y''+2xy'-2y=0;  \quad \{x,\,x^2,\,e^x\}$

\part{f}
$(2x-1)y^{(4)}-4xy'''+(5-2x)y''+4xy'-4y=0;
\quad\{x,\,e^x,\,e^{-x},e^{2x}\}$

\part{g} $xy^{(4)}-y'''-4xy'+4y'=0;  \quad\{1,x^2,\,e^{2x},\,e^{-2x}\}$

\item\label{exer:9.1.7} Find the Wronskian $W$ of a set of three solutions
of $$ y'''+2xy''+e^xy'-y=0, $$ given that $W(0)=2$.

\item\label{exer:9.1.8}
Find the Wronskian $W$ of a set of four solutions of
$$
y^{(4)}+(\tan x)y'''+x^2y''+2xy=0,
$$
given that $W(\pi/4)=K$.

\item\label{exer:9.1.9}
\begin{alist}
\item %(a)
Evaluate the Wronskian $W$ $\{e^x,\,xe^x,\,
x^2e^x\}$. Evaluate $W(0)$.

\item %(b)
Verify that $y_1$, $y_2$, and $y_3$ satisfy
$$
y'''-3y''+3y'-y=0.
\eqno{\rm (A)}
$$

\item %(c)
Use $W(0)$ from \part{a} and Abel's formula to calculate $W(x)$.

\item %(d)
What is the general solution of  (A)?
\end{alist}

\item\label{exer:9.1.10}
Compute the Wronskian of the given set of functions.

\noindent\begin{tabular}[t]{@{}p{168pt}@{}p{168pt}}
\part{a} $\{1,\,e^x,\,e^{-x}\}$
&\part{b} $\{e^x,\, e^x\sin x,\,e^x\cos x\}$
\end{tabular}

\noindent\begin{tabular}[t]{@{}p{168pt}@{}p{168pt}}
\part{c} $\{2,\,x+1,\,x^2+2\}$
&\part{d} $ x,\,x\ln x,\,1/x\}$
\end{tabular}

\begin{tabular}[t]{@{}p{168pt}@{}p{168pt}}
 \part{e} $\dst{\{1,\,x,\,{x^2\over2!},\,
{x^3\over3!}\,,\cdots,\,{x^n\over n!}}\}$&
\part{f} $\{e^x,\,e^{-x},\,x\}$
\end{tabular}

\begin{tabular}[t]{@{}p{168pt}@{}p{168pt}}
\part{g}  $\{e^x/x,\,e^{-x}/x,\,1\}$&
\part{h}  $\{x,\,x^2,\,e^x\}$
\end{tabular}

\begin{tabular}[t]{@{}p{168pt}@{}p{168pt}}
\part{i} $\{x,\,x^3,\,1/x,\,1/x^2\}$&
\part{j}  $\{e^x,\,e^{-x},\,x,\,e^{2x}\}$
\end{tabular}

\part{k} $\{e^{2x},\,e^{-2x},\,1,\,x^2\}$

\item\label{exer:9.1.11}
Suppose $Ly=0$ is normal on $(a,b)$ and $x_0$ is in $(a,b)$. Use
Theorem~\ref{thmtype:9.1.1} to show that $y\equiv0$ is the only solution
of the initial value problem
$$
Ly=0, \quad  y(x_0)=0,\quad y'(x_0)=0,\dots, y^{(n-1)}(x_0)=0,
$$
on $(a,b)$.

\item\label{exer:9.1.12}
Prove:  If $y_1$, $y_2$, \dots, $y_n$ are solutions of $Ly=0$ and the
functions
$$
z_i=\sum^n_{j=1}a_{ij}y_j,\quad 1\le i\le n,
$$
form a fundamental set of solutions of $Ly=0$, then so do $y_1$, $y_2$,
\dots, $y_n$.

\item\label{exer:9.1.13}
Prove:  If
$$
y=c_1y_1+c_2y_2+\cdots+c_ky_k+y_p
$$
is a solution of a linear equation $Ly=F$ for every choice of the
constants $c_1$, $c_2$ ,\dots, $c_k$, then $Ly_i=0$ for $1\le i\le k$.

\item\label{exer:9.1.14}
Suppose $Ly=0$ is normal on $(a,b)$ and let $x_0$ be in $(a,b)$.
For $1\le i\le n$, let $y_i$ be the solution of the initial value
problem
$$
Ly_i=0, \quad  y_i^{(j)} (x_0)=
\left\{\begin{array}{cl}
0,& j\ne i-1,\\ [1\jot]
1,&j=i-1,\end{array}\right. 1\le i\le n,
$$
where $x_0$ is an arbitrary point in $(a,b)$.  Show that any solution
of $Ly=0$ on $(a, b)$, can be written as
$$
y=c_1y_1+c_2y_2+\cdots+c_ny_n,
$$
with $c_j=y^{(j-1)}(x_0)$.

\item\label{exer:9.1.15}
Suppose $\{y_1, y_2,\dots, y_n\}$ is a fundamental set of
solutions of
$$
\dst P_0(x)y^{(n)}+P_1(x)y^{(n-1)}+\cdots+P_n(x)y=0
$$
on $(a,b)$, and let
$$
\begin{array}{rl}
z_1&=a_{11}y_1+a_{12}y_2+\cdots+a_{1n}y_n\\
z_2&=a_{21}y_1+a_{22}y_2+\cdots+a_{2n}y_n\\
\phantom{z_1}&\vdots\phantom{_1y_1+a}\vdots
\phantom{_2y_2+\cdots+a}\vdots\phantom{_ny_n} \phantom{=b}\vdots\\
z_n&=a_{n1}y_1+a_{n2}y_2+\cdots+a_{nn}y_n,
\end{array}
$$
where the $\{a_{ij}\}$ are constants. Show that $\{z_1, z_2,\dots,
z_n\}$ is a fundamental set of solutions  of (A) if and only if the
determinant
$$
\left|\begin{array}{cccc}
a_{11}&a_{12}&\cdots&a_{1n}\\
a_{21}&a_{22}&\cdots&a_{2n}\\
\vdots&\vdots&\ddots&\vdots\\
a_{n1}&a_{n2}&\cdots&a_{nn}\end{array}\right|
$$
is nonzero.\hint{The determinant of a product of  $n\times
n$ matrices equals the product of the determinants.}

\item\label{exer:9.1.16}
Show that $\{y_1,y_2,\dots,y_n\}$ is linearly dependent on $(a,b)$ if
and only if at least one of the functions $y_1$, $y_2$, \dots, $y_n$ can be
written as a linear combination of the others on $(a,b)$.

\exercisetext{Take the following as a hint in
Exercises~\ref{exer:9.1.17}--\ref{exer:9.1.19}:

\color{blue}\it
 By the definition of determinant,
$$
\left|\begin{array}{cccc}
a_{11}&a_{12}&\cdots&a_{1n}\\[2\jot]
a_{21}&a_{22}&\cdots&a_{2n}\\[2\jot]
\vdots&\vdots&\ddots&\vdots\\[2\jot]
a_{n1}&a_{n2}&\cdots&a_{nn}\end{array}\right|
=\sum\pm a_{1i_1}a_{2i_2},\dots,a_{ni_n},
$$
where the sum is over all permutations $(i_1,i_2,\dots,i_n)$ of
$(1,2,\dots, n)$ and the choice of $+$ or $-$ in each term depends
only on the permutation associated with that term.}
\color{black}\rm

\item\label{exer:9.1.17}
Prove:  If
$$
A(u_1,u_2,\dots,u_n)=
\left|\begin{array}{cccc}
a_{11}&a_{12}&\cdots&a_{1n}\\[2\jot]
a_{21}&a_{22}&\cdots&a_{2n}\\[2\jot]
\vdots&\vdots&\ddots&\vdots\\[2\jot]
a_{n-1,1}&a_{n-1,2}&\cdots&a_{n-1,n}\\[2\jot]
u_1&u_2&\cdots&u_n\end{array}\right|,
$$
then
$$
A(u_1+v_1, u_2+v_2,\dots,
  u_n+v_n)=A(u_1,u_2,\dots,u_n)+A(v_1,v_2,\dots,
  v_n).
$$

\item\label{exer:9.1.18}
Let
$$
F=\left|\begin{array}{cccc}
f_{11}&f_{12}&\cdots&f_{1n}\\[2\jot]
f_{21}&f_{22}&\cdots&f_{2n}\\[2\jot]
\vdots&\vdots&\ddots&\vdots\\[2\jot]
f_{n1}&f_{n2}&\cdots&f_{nn}\end{array}\right|,
$$
where $f_{ij}\; (1\le i,\;   j\le n)$ is differentiable.  Show that
$$
F'=F_1+F_2+\cdots+F_n,
$$
where $F_i$ is the determinant obtained by differentiating the
$i$th row of $F$.

\item\label{exer:9.1.19}
Use Exercise~\ref{exer:9.1.18}  to show that if $W$ is the
Wronskian of the
$n$-times differentiable functions $y_1$, $y_2$, \dots, $y_n$, then
$$
W'=
\left|\begin{array}{cccc}
y_1&y_2&\cdots&y_n\\[2\jot]
y'_1&y'_2&\cdots&y'_n\\[2\jot]
\vdots&\vdots&\ddots&\vdots\\[2\jot]
y_1^{(n-2)}&y_2^{(n-2)}&\cdots&y_n^{(n-2)}\\[2\jot]
y_1^{(n)}&y_2^{(n)}&\cdots&y_n^{(n)}
\end{array}\right|.
$$

\item\label{exer:9.1.20}
Use Exercises~\ref{exer:9.1.17} and \ref{exer:9.1.19} to show that if $W$ is
the Wronskian
of  solutions $\{y_1,y_2,\dots,y_n\}$ of the normal equation
$$
P_0(x)y^{(n)}+P_1(x)y^{(n-1)}+\cdots+P_n(x)y=0,
\eqno{\rm (A)}
$$
then $W'=-P_1W/P_0$.
Derive Abel's formula (Eqn.~\eqref{eq:9.1.15}) from this.
\hint{Use (A) to write $y^{(n)}$ in terms of
$y,y',\dots,y^{(n-1)}$.}

\item\label{exer:9.1.21}  Prove Theorem~\ref{thmtype:9.1.6}.

\item\label{exer:9.1.22}  Prove Theorem~\ref{thmtype:9.1.7}.

\item\label{exer:9.1.23}
Show that if the Wronskian of the $n$-times continuously
differentiable functions $\{y_1,y_2,\dots,y_n\}$ has no zeros in
$(a,b)$, then the differential equation obtained by expanding the
determinant
$$
\left|\begin{array}{ccccc}
y&y_1&y_2&\cdots&y_n\\[2\jot]
y'&y'_1&y'_2&\cdots&y'_n\\[2\jot]
\vdots&\vdots&\vdots&\ddots&
 \vdots\\[2\jot]
y^{(n)}&y_{1}^{(n)}&y_2^{(n)}&\cdots&y_n^{(n)}
\end{array}\right|=0,
$$
in cofactors of its first column is normal and has
$\{y_1,y_2,\dots,y_n\}$ as a fundamental set of solutions on $(a,b)$.

\item\label{exer:9.1.24}
Use the method suggested by Exercise~\ref{exer:9.1.23}  to find a linear
homogeneous equation such that the given set of functions
is a fundamental set of solutions on intervals on which the Wronskian
of the set has no zeros.

\begin{tabular}[t]{@{}p{168pt}@{}p{168pt}}
\part{a} $\{x,\,x^2-1,\,x^2+1\}$&
\part{b} $\{e^x,\,e^{-x},\,x\}$
\end{tabular}

\begin{tabular}[t]{@{}p{168pt}@{}p{168pt}}
\part{c}  $\{e^x,\,xe^{-x},\,1\}$&
\part{d}  $\{x,\,x^2,\,e^x\}$
\end{tabular}

\begin{tabular}[t]{@{}p{168pt}@{}p{168pt}}
\part{e} $\{x,\,x^2,\,1/x\}$&
\part{f} $\{x+1,\,e^x,\,e^{3x}\}$
\end{tabular}

\begin{tabular}[t]{@{}p{168pt}@{}p{168pt}}
 \part{g} $\{x,\,x^3,\,1/x,\,1/x^2\}$&
\part{h} $\{x,\,x\ln x,\,1/x,\,x^2\}$
\end{tabular}

\begin{tabular}[t]{@{}p{168pt}@{}p{168pt}}
\part{i} $\{e^x,\,e^{-x},\,x,\,e^{2x}\}$&
\part{j} $\{e^{2x},\,e^{-2x},\,1,\,x^2\}$
\end{tabular}

\end{exerciselist}
